\emph{(i).} The field theory's unit is $i = I_{4}$ (Proposition \claimref{field-theory.second-quantisation}(4)). $B$ is even, so $I_{4}B = BI_{4}$ (Proposition \claimref{field-theory.where-each-lives}(1)), and $(I_{4}B)^{2} = I_{4}^{2}B^{2} = -1$. For $B = t \wedge u = tu$ with $t$ timelike, $u$ spacelike, $t \cdot u = 0$, $I_{4}B$ is the 2-blade of the orthogonal complement of $\mathrm{span}\{t, u\}$ (for $B = e_{0}e_{1}$, $I_{4}B = e_{2}e_{3}$), a spacelike plane; a vector of $B$ anticommutes with $I_{4}$ and with $B$, so it commutes with $I_{4}B$, and a vector $n$ of $B^{*}$ anticommutes with $I_{4}$ and commutes with $B$, so $B^{*}n = -nB^{*}$. Since $\lambda B$ and $I_{4}\theta B$ commute, $\exp(-(\lambda + I_{4}\theta)B/2) = \exp(-\lambda B/2)\exp(-\theta B^{*}/2)$, and the powers of $B^{*}$ cycle $1, B^{*}, -1, -B^{*}$, so the series is $\cos\tfrac{\theta}{2} - B^{*}\sin\tfrac{\theta}{2}$. The check verifies the blade, square, commutation and power-cycle statements exactly for the three coordinate planes $e_{0}e_{k}$, a rotated plane $e_{0}\wedge(\tfrac35e_{1} + \tfrac45e_{2})$ and two boosted planes, including $I_{4}^{k}B^{k} = (B^{*})^{k}$ for $k \le 4$.

\emph{(ii).} With the exact half-angle values, $R(2\pi) = \cos\pi = -1$, $R(4\pi) = 1$ and $R(\pi) = -B^{*}$. Left multiplication by $-1$ sends every even element to its negative; conjugation by $-1$ fixes everything. A spinor value transforms by left action $\psi \mapsto R\psi$ and a bivector by $F \mapsto RF\tilde R$, so a bilinear $\psi F_{0}\tilde\psi$ goes to $R\psi F_{0}\tilde\psi\tilde R$ and returns to itself at $2\pi$; a value with $2s$ spinor factors picks up $(-1)^{2s}$. $R(\pi) = -B^{*}$ and $\tilde R(\pi) = B^{*}$, so conjugation sends $v \mapsto -B^{*}vB^{*}$: for $v$ in $B$, $v$ commutes with $B^{*}$ and the image is $-v(B^{*})^{2} = v$; for $v$ in $B^{*}$ it anticommutes and the image is $-v$. All of this is checked exactly on the eight even basis elements, the six basis bivectors, one spinor bilinear and the vectors of both planes for every sample plane. Points are vectors, and $I_{4}$ anticommutes with vectors, so the factorisation of (i) does not hold for the action on points: there the continuation is the complex rotation of Proposition \claimref{singularity.euclidean-rotor-where-it-lies}(3), realised by no real rotor; on values it is the real rotor of the dual plane, which for the radial boost plane is the angular plane $I\sigma_{r}$ of that proposition. The stored continued value factor of Theorem \claimref{synthesis.boost-evolution-on-generators-is-one-unitary-group}(iv), $-1$ at $\theta = 2\pi$ and $+1$ at $4\pi$ on the canonical spinor component, is the case $s = \tfrac12$, and with the unit $I_{4}$ its rotor $\cos\tfrac{\theta}{2} + I_{4}e_{1}e_{0}\sin\tfrac{\theta}{2}$ is the real rotor $\cos\tfrac{\theta}{2} - e_{2}e_{3}\sin\tfrac{\theta}{2}$. On a tower plane's own modes the imaginary-rapidity rotor is the unitary phase flow of Proposition \claimref{field-theory.boost-modes}; the sign at the full period on values is a separate factor carried by the field's spinor index.

\emph{(iii).} For the completed pair $v = \tfrac12(e_{0} + e_{1})$, $v^{*} = \tfrac12(e_{0} - e_{1})$ of $\Cl(1,1)$, with $v^{*}v$ the occupied projector (Proposition \claimref{hawking.modular-structure-completed-pair}), the Gibbs state $\rho = (vv^{*} + x\,v^{*}v)/(1 + x)$, $\omega(X) = 2\langle\rho X\rangle_{0}$, has $\omega(1) = 1$, $\omega(v^{*}v) = x/(1+x)$ and $\omega(v^{*}v) = x\,\omega(vv^{*})$, the KMS relation for the phase flow $a \mapsto e^{-iw\lambda}a$ at period $2\pi$ on the monomials; with $\omega(vv^{*}) = 1 - \omega(v^{*}v)$ it reads $n = x(1 - n)$. For a CCR mode the Gibbs occupation $n = x/(1 - x)$ satisfies $n = x(n + 1)$, the same KMS relation with $\omega(bb^{*}) = \omega(b^{*}b) + 1$. Both are exact at $x = \tfrac13, \tfrac27, \tfrac5{11}$. The KMS condition fixes the state once the algebra is given; it does not choose between the two algebras. Theorem \claimref{synthesis.local-boost-temperature-is-local-regularity}(iii) states (R$_{\mathrm{loc}}$) on CAR windows, so the statistics there are an input.

\emph{(iv).} Imported: D. Guido and R. Longo, \emph{An algebraic spin and statistics theorem}, Commun. Math. Phys. (1995), arXiv funct-an/9406005, the field-net theorem of Section 2, Thm 2.11 (the source's numbering also uses 2.11 for an equation; Fact \claimref{synthesis.imported-free-field-spin-statistics-and-guido-longo-theorem}(iii), Commun. Math. Phys. 172). Its hypotheses are a modular covariant net of field algebras (the wedge modular groups act as the wedge boosts, rescaled to rapidity $-2\pi t$), the Reeh-Schlieder property and normal Bose/Fermi commutation relations with statistics operator $\Gamma$; positive energy is derived, not assumed, and the conclusion is $\Gamma = U(2\pi)$. The proof there does not assume the spin-statistics relation. The Bisognano-Wichmann theorem (memory: J. Math. Phys. 16 (1975) 985; 17 (1976) 303), which supplies modular covariance for Wightman fields, is not used here, since its derivation starts from Wightman fields with the normal commutation relations and would make the argument circular. $U(2\pi)$ acts as $(-1)^{2s}$ on a spin-$s$ value by representation theory alone (the $2\pi$ rotation of $\Spin$ is $-1$ on the spinor representation), so under modular covariance $\Gamma = (-1)^{2s}$ on the spin-$s$ fields; (i) and (ii) are not needed for this. They re-express it: on values $U(2\pi)$ is the continued wedge boost over one modular period, and the half-period rotor $R(\pi)$ is the turn by $\pi$ of the dual plane, the rotation that appears with the total inversion in the Bisognano-Wichmann form of the wedge's modular conjugation (memory). On vector arguments the complexified boost at half the period acts as $-1$ on the vectors of $B$ and $+1$ on those of $B^{*}$, which is the dual turn ($+1$ on $B$, $-1$ on $B^{*}$) composed with the total inversion. This is standard content (on the two chiral spinors the boost generator is $\pm$ the unit times the rotation generator) re-derived in the formalism; the series-specific part is only that with the unit $I_{4}$ the continuation on values is a real rotor. Modular covariance with the geometric action on values is stronger than (R$_{\mathrm{loc}}$) as stored, and is not a premise of the store; the spin-statistics input can be traded for it, not derived from stored premises. For the attached free fields the statistics follow from locality with a positive norm, the free-field route (\claimref{synthesis.locality-of-the-attached-free-fields-fixes-their-statistics}); the Jost-point import remains for the interacting case.
